\originalchapter{太阳内部、太阳风与变星的物理时钟}
本章支撑8.901第12、19、22次课。第19次太阳课在官方课历中署Justin Kasper客座；没有取得其课堂全文，因此不把以下原创推导冒充客座原讲义。太阳的热输运、外层流和恒星脉动提供了互补的时间与密度信息。

\originalsection{太阳模型的三类检验}
球对称结构把质量、半径、光度与组成连在一起；这些全球数足以约束，但不唯一决定每层温度和密度。核心H燃烧供能，光在高光学深度区扩散，外层不透明度和EOS变化可使辐射梯度超过对流阈值，产生对流包层。星体温度从核心向外下降，而光谱给$T_{\rm eff}$，不能把表面几千度直接作为核燃烧温度。

随机走光子的平均自由程$\ell$远小于$R$，走出长度$R$需要约$(R/\ell)^2$次步进，所以扩散时间量级$t_{\rm diff}\sim R^2/(c\ell)=\tau R/c$；几何、密度不均和输运系数使系数改变。中微子极弱相互作用可近自由逃逸，因而同一次核反应的信息不与光同时离开星体。

太阳中微子测核反应及味变换，日震测声速、分层和流，表面光谱测组成。声波局部$c_s^2=\Gamma_1P/\rho$，频率谱取决于声腔边界、温度和组成；从有限本征频率恢复内部是带模型与误差的逆问题，不等于直接看见每一层。PS5.4用双味振荡展示相位差怎样改变检测到的电子中微子比例，不把双味真空模型称为全部太阳MSW计算。

\originalsection{Parker太阳风的临界点}
静态高温外大气并非唯一可能。取稳态、球对称、等温气体，等温声速$c_T$恒定，连续$\dot M=4\pi r^2\rho v$，动量
\[
 v\frac{\dd v}{\dd r}=-c_T^2\frac{\dd\ln\rho}{\dd r}-\frac{GM}{r^2}.
\]
连续式$\dd\ln\rho/\dd r=-2/r-v'/v$代入：
\[
 \left(v-\frac{c_T^2}{v}\right)\frac{\dd v}{\dd r}
 =\frac{2c_T^2}{r}-\frac{GM}{r^2}.
\]
若要平滑穿过$v=c_T$，右侧也须为零，临界半径$r_s=GM/(2c_T^2)$。积分并选择临界点常数，得
\[
 \frac{v^2}{c_T^2}-\ln\frac{v^2}{c_T^2}
 =4\ln\frac r{r_s}+\frac{4r_s}{r}-3 .
\]
同一隐式关系有亚声、超声支，边界和正向质量流选择跨临界的外流；仅写$P=\rho c_T^2$不能唯一选择解。真实太阳风有加热、磁场、非等温与多组分，以上揭示临界点结构，不是精确观测拟合。热星辐射驱动风与这里压力驱动风也不同，PS3.1和PS6.4分别分析指定辐射作用。

\originalsection{脉动周期连接平均密度}
径向振动中回位重力与压力变化的典型频率$\omega^2\sim GM/R^3$，所以$P_{\rm puls}\sim Q\sqrt{R^3/(GM)}$。无量纲$Q$由密度剖面、振动阶数、EOS与非绝热热交换决定，不能从量纲给一个精确普适数。更稀疏、半径大的星振动慢，这是Cepheid类变星可关联光度与周期的结构基础。

部分电离层的$\kappa$机制要求适当相位：压缩时不透明度提高，辐射暂被阻留并增加压强，在有利阶段向振动作功；随后膨胀释放能量。只有热调整时标与振动周期匹配、驱动超过其他层阻尼时，净周期功为正。不是“有一处不透明度增大”就必定出现持续振动。

若不稳定带温度范围较窄，用$L\propto R^2T_{\rm eff}^4$将半径连到光度，再用周期--密度关系、质量光度关系得到周期与光度的相关。校准零点仍需独立距离；金属丰度、消光及不同脉动模式改变关系。多波段或组合星等可减少消光，但不会自动消除全部系统误差。

其他变量机制不同：食双星是几何遮挡，吸积盘不稳定是输运／热反馈，白矮星表面nova是积累氢后热失控，中子星X射线暴可能是表面燃烧。它们不应都被同一恒星径向脉动公式解释。
